\section{Catenary schemes}
\label{section-catenary}

\noindent
Recall that a topological space $X$ is called {\it catenary} if
for every pair of irreducible closed subsets $T \subset T'$
there exist a maximal chain of irreducible closed subsets
$$
T = T_0 \subset T_1 \subset \ldots \subset T_e = T'
$$
and every such chain has the same length. See
Topology, Definition \ref{topology-definition-catenary}.

\begin{definition}
\label{definition-catenary}
Let $S$ be a scheme. We say $S$ is {\it catenary} if the
underlying topological space of $S$ is catenary.
\end{definition}

\noindent
Recall that a ring $A$ is called {\it catenary} if
for any pair of prime ideals $\mathfrak p \subset \mathfrak q$
there exists a maximal chain of primes
$$
\mathfrak p =
\mathfrak p_0 \subset \ldots \subset \mathfrak p_e
= \mathfrak q
$$
and all of these have the same length. See
Algebra, Definition \ref{algebra-definition-catenary}.

\begin{lemma}
\label{lemma-catenary-local}
Let $S$ be a scheme. The following are equivalent
\begin{enumerate}
\item $S$ is catenary,
\item there exists an open covering of $S$ all of whose members are
catenary schemes,
\item for every affine open $\Spec(R) = U \subset S$ the ring
$R$ is catenary, and
\item there exists an affine open covering $S = \bigcup U_i$ such
that each $U_i$ is the spectrum of a catenary ring.
\end{enumerate}
Moreover, in this case any locally closed subscheme of $S$ is catenary
as well.
\end{lemma}

\begin{proof}
Combine Topology, Lemma \ref{topology-lemma-catenary}, and
Algebra, Lemma \ref{algebra-lemma-catenary}.
\end{proof}

\begin{lemma}
\label{lemma-catenary-dimension-function}
Let $S$ be a locally Noetherian scheme.
The following are equivalent:
\begin{enumerate}
\item $S$ is catenary, and
\item locally in the Zariski topology there exists a dimension function
on $S$ (see Topology, Definition \ref{topology-definition-dimension-function}).
\end{enumerate}
\end{lemma}

\begin{proof}
This follows from
Topology, Lemmas
\ref{topology-lemma-catenary},
\ref{topology-lemma-dimension-function-catenary}, and
\ref{topology-lemma-locally-dimension-function},
Schemes, Lemma \ref{schemes-lemma-scheme-sober}
and finally Lemma \ref{lemma-Noetherian-topology}.
\end{proof}

\noindent
It turns out that a scheme is catenary if and only if its local
rings are catenary.

\begin{lemma}
\label{lemma-catenary-local-rings-catenary}
Let $X$ be a scheme. The following are equivalent
\begin{enumerate}
\item $X$ is catenary, and
\item for any $x \in X$ the local ring $\mathcal{O}_{X, x}$ is
catenary.
\end{enumerate}
\end{lemma}

\begin{proof}
Assume $X$ is catenary. Let $x \in X$. By Lemma \ref{lemma-catenary-local}
we may replace $X$ by an affine open neighbourhood of $x$, and
then $\Gamma(X, \mathcal{O}_X)$ is a catenary ring. By
Algebra, Lemma \ref{algebra-lemma-localization-catenary} any
localization of a catenary ring is
catenary. Whence $\mathcal{O}_{X, x}$ is catenary.

\medskip\noindent
Conversely assume all local rings of $X$ are catenary.
Let $Y \subset Y'$ be an inclusion of irreducible closed
subsets of $X$. Let $\xi \in Y$ be the generic point.
Let $\mathfrak p \subset \mathcal{O}_{X, \xi}$ be the prime
corresponding to the generic point of $Y'$, see
Schemes, Lemma \ref{schemes-lemma-specialize-points}. By that same
lemma the irreducible closed subsets of $X$ in between $Y$ and $Y'$
correspond to primes $\mathfrak q \subset \mathcal{O}_{X, \xi}$
with $\mathfrak p \subset \mathfrak q \subset \mathfrak m_{\xi}$.
Hence we see all maximal chains of these are finite and have the
same length as $\mathcal{O}_{X, \xi}$ is a catenary ring.
\end{proof}






